\chapter{Parathyroid glands}
\index{Parathyroid glands}

\medskip
\noindent\textit{Educational information; seek professional advice for individual care.}

\section*{What they do}
The parathyroid glands are usually four small glands near the thyroid. They release parathyroid hormone, which helps regulate calcium and phosphate levels and supports nerve, muscle, and bone function. Problems may cause too much or too little hormone.

\section*{Overactive glands}
Primary hyperparathyroidism raises blood calcium and may cause no symptoms. Possible features include thirst, frequent urination, constipation, nausea, fatigue, weakness, low mood, bone pain, kidney stones, or reduced bone density. Causes include a benign overactive gland, enlargement of several glands, inherited conditions, or rarely cancer.

\section*{Assessment}
Diagnosis uses repeated blood tests for calcium and parathyroid hormone, with kidney function, phosphate, vitamin D, urine calcium, bone-density testing, or imaging when indicated. Imaging helps plan surgery; it does not by itself diagnose the condition.

\section*{Treatment}
Some people need monitoring, while others benefit from surgery to remove the overactive gland or glands. Treatment of vitamin D deficiency, kidney disease, medicines, and fluid balance is individualized. Do not take large calcium or vitamin D supplements without advice.

\section*{When to seek urgent care}
Seek urgent help for confusion, severe weakness, repeated vomiting, severe dehydration, fainting, a severe irregular heartbeat, or sudden worsening illness. Symptoms of low calcium such as muscle spasms, tingling around the mouth, or seizures also need urgent assessment.

\section*{Sources}
\begin{itemize}
\item National Institute of Diabetes and Digestive and Kidney Diseases: \url{https://www.niddk.nih.gov/health-information/endocrine-diseases/primary-hyperparathyroidism}
\end{itemize}
